\documentclass[10pt,a4paper]{amsart}
\usepackage[margin=19mm]{geometry}
\usepackage{amsmath,amssymb,mathtools}
\usepackage[scr=rsfs,cal=euler]{mathalfa}
\usepackage{xfrac}
\usepackage[unicode,hidelinks,bookmarks=false]{hyperref}
\newcommand{\ie}{i.e.\ }
\newcommand{\cf}{cf.\ }
\newcommand{\resp}{respectively\ }
\newcommand{\Subsection}{\subsection}
\providecommand{\nobreakdash}{\nobreak\hbox{-}}
\setlength{\emergencystretch}{2em}
\multlinegap0pt
\usepackage{bm}
\usepackage{bbm}
\usepackage{dsfont}
\usepackage{xspace}
\usepackage{cancel}
\usepackage{mathtools}
\usepackage{amsthm}
\makeatletter
\@ifundefined{proposition}{\newtheorem{proposition}{Proposition}}{}
\makeatother
\renewcommand{\geq}{\geqslant}
\renewcommand{\leq}{\leqslant}
\title[Multiplicative recurrence of M\"obius transformations]{Multiplicative recurrence of M\"obius transformations}
\author{Sun-Kai Leung, Christian T\'afula}
\date{Source: arXiv:2409.12936v3 (CC BY 4.0)}
\begin{document}
\maketitle

\noindent\emph{Source and licence.} Multiplicative recurrence of M\"obius transformations. Version arXiv:2409.12936v3 (CC BY 4.0), distributed under the Creative Commons Attribution 4.0 International licence, as recorded in the arXiv licence metadata for this exact version and shown on the abstract page https://arxiv.org/abs/2409.12936v3. Full rights notice: licenses/arxiv-2409.12936-CC-BY-4.0.txt.\par\smallskip

\noindent\textit{Editorial scope.} Counted unit: the Proposition whose source label is \texttt{prop:a=b(b-1)} in the article (source0.tex lines 695--698, source label \texttt{prop:a=b(b-1)}), delivered together with the article's own complete original proof of it (source0.tex lines 700--747, one \texttt{proof} environment of 48 lines). No mathematics is written, repaired or re-derived: statement and proof are the article's own text line for line. Required context retained uncounted: none. Every definition, display and label of those lines is retained unchanged. Omitted from this package and not redistributed: 1--694, 699--699, 748--918. Nothing inside a retained excerpt is omitted or altered: every retained body is delivered exactly as the article's lines are written, so no omission receipt of any kind accompanies this package. Citations: the retained excerpts carry no citation command at all, so no bibliography entry of the article is transcribed and no reference is redistributed. Reference adaptation: none was needed and none was made; every reference inside the delivered unit resolves inside it, so no unresolved reference or citation is produced. Printed slips kept literal: no printed word, symbol, hypothesis or number of the counted statement or of its proof was altered. Layout and class: the article's own class is \texttt{amsart} with packages including amssymb, amsmath, enumerate, amsfonts, amsthm, mathrsfs, url, bm, mathtools, xcolor, hyperref, color, geometry, doi; the wrapper is the lane's pinned \texttt{amsart} editorial wrapper at 10pt on A4, which restates the theorem-like environments the retained text needs and adds the article's own macro definitions that the retained text needs (\texttt{\textbackslash geq}, \texttt{\textbackslash leq}), with extra wrapper packages (bm, bbm, dsfont, xspace, cancel, mathtools). The delivered numbering is the unit's own. Assets: no retained span contains a figure environment, a graphics-inclusion command, a PDF-page inclusion, a table or a TikZ picture; no image file is copied into this package, the package declares no asset dependency and no image file is redistributed with it. Licence: version arXiv:2409.12936v3 of this article is distributed under the Creative Commons Attribution 4.0 International licence, as recorded in the arXiv licence metadata for this exact version and shown on the abstract page https://arxiv.org/abs/2409.12936v3. A full rights notice is delivered at licenses/arxiv-2409.12936-CC-BY-4.0.txt. Characters: every retained excerpt is pure ASCII, so the delivered engine, XeLaTeX with its default Latin Modern text font, reports no missing character.
\begin{proposition} \label{prop:a=b(b-1)}
Let $a, b \in \mathbb{N}$ with $a=b(b-1).$ Then
for any integer $k \geq 2,$ there exists an integer $H_k \geq 1$ depending on $a,b$ such that, for each integer $n\geq 0$, the induced subgraph $G_{a,b,a,b-1}[\{n+1, \ldots, n+H_k\}]$ contains a $k$-clique. In particular, we have $\omega(G_{a,b,a,b-1})=\infty.$
\end{proposition}

\begin{proof}
We proceed by induction on $k$. When \( k = 2 \), one can simply take \( H_2 = a + b \), since every interval of length \( H_2 \) contains a pair of consecutive integers of the form \( an + (b - 1) \), \( an + b \).
 Suppose that the statement holds for some integer $k \geq 2.$ Let 
\begin{align*}
M_k:=\prod_{1 \leq h \leq H_k}(bh+1) \quad \text{and} \quad r_k:=\frac{1}{b}(M_k-1).
\end{align*}
Then by the induction hypothesis, there exist integers $1 \leq h_1, \ldots, h_k \leq H_k$ such that the set of vertices $\{(b-1)r_k+h_1, \ldots, (b-1)r_k+h_k\}$ forms a $k$-clique in $G_{a,b,a,b-1}.$ In other words, for any integers $1 \leq i < j \leq k,$ we have
\begin{align*}
\frac{(b-1)r_k+h_i}{(b-1)r_k+h_j}=\frac{am_{ij}+b}{am_{ij}+(b-1)}
\end{align*}
for some integers $m_{ij} \geq 1.$ Since 
\begin{align*}
\frac{am_{ij}+b}{am_{ij}+(b-1)}
\end{align*}
is a reduced fraction, we have
\begin{gather*}
(b-1)r_k+h_i=g_{ij} (am_{ij}+b) \quad \text{and} \quad
(b-1)r_k+h_j= g_{ij}(am_{ij}+(b-1)),
\end{gather*}
where $g_{ij}:=\gcd((b-1)r_k+h_i, (b-1)r_k+h_j)$. It follows that $g_{ij}= h_i-h_j \leq H_{k}$ by considering their difference.
%\[ am+b = \frac{(b-1)r_k+h_i}{g_{ij}} = \frac{(b-1)r_k+h_j}{g_{ij}} + \frac{h_i-h_j}{g_{ij}} = am+(b-1) + \frac{h_i-h_j}{g_{ij}}. \] 
Let
\begin{align*}
n_{\ell} := b (M_k+H_k)! {\ell}+ r_k %\label{eq:n_t}
\end{align*}
for integers ${\ell} \geq 1.$ Then by the assumption $a=b(b-1)$, we have
\begin{align}
\frac{an_{\ell}+bh_i}{an_{\ell}+bh_j} &= \frac{an_{\ell}+b(g_{ij} (am_{ij}+b)-(b-1)r_k)}{an_{\ell}+b(g_{ij}(am_{ij}+(b-1))-(b-1)r_k)} \nonumber \\
&=  \frac{a  b (M_k+H_k)! {\ell}+bg_{ij}(am_{ij}+b)}{a  b (M_k+H_k)! {\ell}+bg_{ij} (am_{ij}+(b-1))}.  \label{eq:crazyfrac}
\end{align}
Since $g_{ij} \leq H_k < M_k+H_k,$ we have $g_{ij} \mid (M_k+H_k)!,$ so that the fraction in (\ref{eq:crazyfrac}) can be reduced to
\begin{equation*}
\frac{an_{\ell}+bh_i}{an_{\ell}+bh_j} = \frac{a \left( \frac{(M_k+H_k)!}{g_{ij}}{\ell}+m_{ij} \right)+b}{a \left( \frac{(M_k+H_k)!}{g_{ij}}{\ell}+m_{ij} \right)+(b-1)} \in R_{a,b,a,b-1}.
\end{equation*}
Therefore, the set of vertices $\{ an_{\ell}+bh_1, \ldots, an_{\ell}+bh_k \}$ forms a new $k$-clique in $G_{a,b,a,b-1}$ for each integer ${\ell} \geq 1$. 

Let ${\ell} \geq 1$ be a fixed integer. We claim that $an_{\ell}+(b-1)$ is adjacent to $an_{\ell}+bh_j$ for each integer $1 \leq j \leq k,$ so that they form a $(k+1)$-clique in $G_{a,b,a,b-1}.$ By our choice of $M_k$, $r_k$, we have $br_k+1\equiv 0 \pmod{bh+1}$ for every $0\leq h\leq H_k$, or equivalently $r_k\equiv h\pmod{bh+1}$. Hence, for each $1\leq h\leq H_k$ there exists an integer $m_h \geq 1$ such that 
\begin{align*}
n_{\ell} = b(M_k+H_k)! {\ell} +r_k = (bh+1)m_h +h.
\end{align*}
Since $a+b = b(b-1)+b = b^2$, this implies that
\begin{align}
\frac{an_{\ell}+b(h+1)}{an_{\ell}+(b-1)} &= \frac{a(bh+1)m_h+ (ah+bh+b)}{a(bh+1)m_h+(ah+b-1)} \nonumber \\
&= \frac{am_h+b}{am_h+(b-1)} \in R_{a,b,a,b-1}. \label{eq:red2}
\end{align}
Therefore, the induced subgraph $G_{a,b,a,b-1}[\{an_{\ell}+(b-1), \ldots, an_{\ell}+b(H_k+1) \}]$ contains a $(k+1)$-clique for each integer ${\ell} \geq 1$. Since $r_k+b(H_k+1) \leq ab(M_k+H_k)!, $ we can take $H_{k+1}=2ab(M_k+H_k)!,$ so that
 the induced subgraph $G_{a,b,a,b-1}[\{n, n+1, \ldots, n+H_{k+1}\}]$ contains a $(k+1)$-clique for all integers $n \geq 1$ as desired, and the proposition follows.
\end{proof}

\end{document}
